% Auto-generated from appendix.aux — do not edit by hand
\newlabel{app:roadmap}{{}{2}{Roadmap: Appendix Sections vs.\ Main Paper}{section*.3}{}}
\newlabel{app:roadmap@cref}{{}{[1][2][]2}}
\newlabel{tab:app-roadmap}{{.1}{2}{Where to look when checking a main-paper claim}{table.caption.4}{}}
\newlabel{tab:app-roadmap@cref}{{[table][1][2147483647,0].1}{[1][2][]2}}
\newlabel{app:related-work-comparison}{{A}{2}{Extended Related Work Comparison}{appendix.A}{}}
\newlabel{app:related-work-comparison@cref}{{[appendix][1][2147483647]A}{[1][2][]2}}
\newlabel{app:methodology}{{B}{2}{Bug Pool Construction and Methodology}{appendix.B}{}}
\newlabel{app:methodology@cref}{{[appendix][2][2147483647]B}{[1][2][]2}}
\newlabel{app:method-pool}{{B.1}{2}{Pool Construction}{subsection.B.1}{}}
\newlabel{app:method-pool@cref}{{[subappendix][1][2147483647,2]B.1}{[1][2][]2}}
\newlabel{tab:dataset-construction-flow}{{B.1}{2}{Separation of the three experimental datasets. The catalog study and the detector benchmark apply different filters to the same evidence corpus}{table.caption.5}{}}
\newlabel{tab:dataset-construction-flow@cref}{{[table][1][2147483647,2]B.1}{[1][2][]2}}
\newlabel{app:method-silent}{{B.2}{2}{Silent-vs-Loud Filtering}{subsection.B.2}{}}
\newlabel{app:method-silent@cref}{{[subappendix][2][2147483647,2]B.2}{[1][2][]2}}
\newlabel{tab:silent-evidence}{{B.2}{3}{Distribution of the 4-tier silent-evidence rating on the 392-record corpus. L1 provides physical buggy/fixed evidence; L2--L3 are textual evidence levels; L4 records have sparse evidence and are retained only as a review queue. Source: \texttt {silent\_evidence\_392.json}}{table.caption.6}{}}
\newlabel{tab:silent-evidence@cref}{{[table][2][2147483647,2]B.2}{[1][2][]3}}
\newlabel{app:method-taxonomy}{{B.3}{3}{Descriptive Subsystem Coding}{subsection.B.3}{}}
\newlabel{app:method-taxonomy@cref}{{[subappendix][3][2147483647,2]B.3}{[1][3][]3}}
\newlabel{app:method-irr}{{B.4}{3}{Inter-Rater Reliability}{subsection.B.4}{}}
\newlabel{app:method-irr@cref}{{[subappendix][4][2147483647,2]B.4}{[1][3][]3}}
\newlabel{tab:irr}{{B.3}{3}{Inter-rater reliability across the two protocols. Protocol A measures re-annotation agreement on 32 doubly-and-independently annotated bugs; Protocol B uses an independent LLM rerater on a 50-bug stratified sample to compute Cohen's $\kappa $. Source: \texttt {irr\_50\_results.json}}{table.caption.9}{}}
\newlabel{tab:irr@cref}{{[table][3][2147483647,2]B.3}{[1][3][]3}}
\newlabel{app:method-limitation}{{B.4}{3}{Label-boundary notes}{section*.10}{}}
\newlabel{app:method-limitation@cref}{{[subappendix][4][2147483647,2]B.4}{[1][3][]3}}
\newlabel{app:method-reproducibility}{{B.5}{3}{From the Evidence Corpus to Real-SE}{subsection.B.5}{}}
\newlabel{app:method-reproducibility@cref}{{[subappendix][5][2147483647,2]B.5}{[1][3][]3}}
\newlabel{app:bug-distribution}{{C}{4}{Fine-Grained Distribution of the 392-Record Evidence Corpus}{appendix.C}{}}
\newlabel{app:bug-distribution@cref}{{[appendix][3][2147483647]C}{[1][4][]4}}
\newlabel{fig:bug-distribution-detail}{{C.1}{4}{Fine-grained distribution of the 392-record evidence corpus. \textit {(a)} By descriptive subsystem label: the background color (sand) marks the corpus, while the foreground (teal) marks historical pre-filter diversity targets used during Real-SE construction; \texttt {s/p} denotes target/pool. \textit {(b)} By dominant training parallelism dimension (None = no obvious dimension dependency; Combo = multi-dimension composite)}{figure.caption.11}{}}
\newlabel{fig:bug-distribution-detail@cref}{{[figure][1][2147483647,3]C.1}{[1][4][]4}}
\newlabel{app:pattern-catalog}{{D}{4}{Template Induction Protocol and Results}{appendix.D}{}}
\newlabel{app:pattern-catalog@cref}{{[appendix][4][2147483647]D}{[1][4][]4}}
\newlabel{tab:template-induction}{{D.1}{4}{Frozen snapshot used in the template-induction study. Schema coverage counts cases matched to a template in that snapshot; joint coverage additionally requires topology and validity guards to be instantiable. Counts describe this experiment, not a fixed system-wide catalog size}{table.caption.13}{}}
\newlabel{tab:template-induction@cref}{{[table][1][2147483647,4]D.1}{[1][4][]4}}
\newlabel{fig:template-induction}{{D.1}{5}{Induction on 186 development cases (200 orderings; 5--95\% bands): coverage ${\approx }0.86$ while templates grow 27$\to $35 (core transfers; tail unsaturated). Early/late points use seed vs.\ expanded catalogs}{figure.caption.16}{}}
\newlabel{fig:template-induction@cref}{{[figure][1][2147483647,4]D.1}{[1][4][]5}}
\newlabel{tab:catalog-direct-details}{{D.2}{5}{Complete matched Catalog ablation. Yield is L4/L2; deployment counts refer to the curated entries produced by each arm}{table.caption.17}{}}
\newlabel{tab:catalog-direct-details@cref}{{[table][2][2147483647,4]D.2}{[1][5][]5}}
\newlabel{app:algorithm}{{E}{5}{Invariant Miner: Complete Algorithm Specification}{appendix.E}{}}
\newlabel{app:algorithm@cref}{{[appendix][5][2147483647]E}{[1][5][]5}}
\newlabel{app:case_study}{{E}{5}{Invariant Miner: Complete Algorithm Specification}{appendix.E}{}}
\newlabel{app:case_study@cref}{{[appendix][5][2147483647]E}{[1][5][]5}}
\newlabel{fig:mining-funnel}{{E.1}{6}{Per-stage candidate counts of the four-stage mining funnel (aggregated across 4 frameworks). F1--F4 are the accept/reject gates above; Deploy is the deployment library after manual cross-framework curation. Callouts match the main-paper verification-funnel figure: $114/400$ clean FP without Stage~3; Skip-F4 $=0/6{,}922$; full funnel $=0/764$. \emph {Note}: F1--F4 here denote funnel levels and are a different axis from the evidence tiers L1--L4 of \S \ref {app:method-silent}}{figure.caption.19}{}}
\newlabel{fig:mining-funnel@cref}{{[figure][1][2147483647,5]E.1}{[1][5][]6}}
\newlabel{app:conf-screening}{{E.2}{6}{Acceptance Rule and Confidence Screening}{subsection.E.2}{}}
\newlabel{app:conf-screening@cref}{{[subappendix][2][2147483647,5]E.2}{[1][6][]6}}
\newlabel{app:io_examples}{{E.3}{6}{Concrete Input/Output Specifications}{subsection.E.3}{}}
\newlabel{app:io_examples@cref}{{[subappendix][3][2147483647,5]E.3}{[1][6][]6}}
\newlabel{app:provenance}{{E.4}{6}{Mining Provenance Audit Protocol}{subsection.E.4}{}}
\newlabel{app:provenance@cref}{{[subappendix][4][2147483647,5]E.4}{[1][6][]6}}
\newlabel{app:sql}{{F}{6}{SQL Query Specification}{appendix.F}{}}
\newlabel{app:sql@cref}{{[appendix][6][2147483647]F}{[1][6][]6}}
\newlabel{tab:json_schema}{{F.1}{7}{JSON Data Field Specification}{table.caption.20}{}}
\newlabel{tab:json_schema@cref}{{[table][1][2147483647,6]F.1}{[1][7][]7}}
\newlabel{app:diagnosis}{{G}{8}{Structured Diagnostic Reports}{appendix.G}{}}
\newlabel{app:diagnosis@cref}{{[appendix][7][2147483647]G}{[1][8][]8}}
\newlabel{app:production-cases}{{H}{8}{Production Case Studies}{appendix.H}{}}
\newlabel{app:production-cases@cref}{{[appendix][8][2147483647]H}{[1][8][]8}}
\newlabel{fig:case-studies}{{H.1}{8}{Diagnostic signals for two production silent-error cases (SFT tokenizer; Megatron Muon+PP+MTP routing). Case~1: losses overlap 2{,}000 steps while SWE-bench diverges 15.7 points; the rule fires at step~1}{figure.caption.21}{}}
\newlabel{fig:case-studies@cref}{{[figure][1][2147483647,8]H.1}{[1][8][]8}}
\newlabel{app:extended_data}{{I}{8}{Extended Experimental Data}{appendix.I}{}}
\newlabel{app:extended_data@cref}{{[appendix][9][2147483647]I}{[1][8][]8}}
\newlabel{tab:detection-results}{{I.1}{9}{Per-case buggy detection results on Real-SE (matching the main-text \emph {Detection Effectiveness and Precision} totals: TrainAudit \NumRealSEDet /\NumRealSE , TrainCheck \NumTCDet /\NumBaselineEval , Na\"ive 0/\NumBaselineEval ). TrainAudit uses native commit replay, whereas the baseline columns report the separate case-aligned surrogate study; the columns are not a same-harness ranking. \textbf {Case ID} uses ``framework abbreviation + fixed commit short hash''; cases without a paired commit ID keep their original identifiers. \textbf {Tier}: A0 = PyTorch public-interface layer; A1 = framework adapter layer. $\checkmark $ = detected; $\times $ = missed. The TrainCheck column shows in parentheses \textit {TC violations / total checks}. LC1/O-003 is the observability-boundary case (subthreshold drift): TrainAudit records the expected miss, while baseline verdicts remain unavailable pending re-run}{table.caption.24}{}}
\newlabel{tab:detection-results@cref}{{[table][1][2147483647,9]I.1}{[1][8][]9}}
\newlabel{fig:tier-coverage}{{I.1}{9}{Cumulative schema-tier coverage of the 392-record evidence corpus. S0 covers 28.3\%; S0--S5/6 cover 77.8\%}{figure.caption.25}{}}
\newlabel{fig:tier-coverage@cref}{{[figure][1][2147483647,9]I.1}{[1][9][]9}}
\newlabel{app:scale-axes}{{I.3}{9}{Query-Path Scaling Microbenchmarks}{subsection.I.3}{}}
\newlabel{app:scale-axes@cref}{{[subappendix][3][2147483647,9]I.3}{[1][9][]9}}
\newlabel{app:amortization}{{I.4}{10}{Amortized Full-Snapshot Overhead}{subsection.I.4}{}}
\newlabel{app:amortization@cref}{{[subappendix][4][2147483647,9]I.4}{[1][9][]10}}
\newlabel{tab:full-snapshot-microbenchmark}{{I.2}{10}{Full-snapshot diagnostic microbenchmark at 1.2B/H20. Values are mean iteration times when one complete snapshot is materialized per step}{table.caption.26}{}}
\newlabel{tab:full-snapshot-microbenchmark@cref}{{[table][2][2147483647,9]I.2}{[1][10][]10}}
\newlabel{fig:amortization}{{I.2}{10}{Amortized full-snapshot overhead vs.\ dump period $K$ (1.2B/H20). The optimized collector meets any budget at a ${\sim }7.7\times $ shorter period}{figure.caption.27}{}}
\newlabel{fig:amortization@cref}{{[figure][2][2147483647,9]I.2}{[1][10][]10}}
\newlabel{app:motivation-figures}{{I.5}{10}{Verified Constraint Library and Hookpoint Deployment}{subsection.I.5}{}}
\newlabel{app:motivation-figures@cref}{{[subappendix][5][2147483647,9]I.5}{[1][10][]10}}
\newlabel{fig:training-panorama}{{I.3}{11}{Panorama of hybrid-parallel LLM training and where silent errors strike: four parallelism dimensions and five highlighted fault loci}{figure.I.3}{}}
\newlabel{fig:training-panorama@cref}{{[figure][3][2147483647,9]I.3}{[1][11][]11}}
\newlabel{tab:realse-class-coverage}{{I.3}{11}{Real-SE diversity across descriptive subsystem labels (representative cases only)}{table.I.3}{}}
\newlabel{tab:realse-class-coverage@cref}{{[table][3][2147483647,9]I.3}{[1][11][]11}}
